% Mục 3.1 (bản rút gọn). Rút từ sections/impl_lock.tex. Nguồn: include/lock.h, backoff.hpp, platform.hpp;
% lệnh AArch64: c++ -O2 -S (Apple clang 21).
Mọi lock trong \texttt{include/lock.h} có đúng \texttt{acquire()} và \texttt{release()} không tham số như
khung của đề. Lấy lock dùng acquire, nhả lock dùng release. Cặp này làm mọi phép ghi trong critical
section trước nhìn thấy được ở chủ lock sau. Trạng thái lock được khai báo \texttt{alignas(kCacheLine)},
nên mỗi lock chiếm riêng một cache line và không có false sharing, tức hai biến độc lập nằm chung một
line và làm mất hiệu lực line của nhau.

TAS lặp \texttt{exchange(true, acquire)} (lệnh \texttt{swpab}) tới khi đọc được \texttt{false}, và nhả
lock bằng \texttt{store(false, release)} (\texttt{stlrb}). TTAS chờ bằng \texttt{load(relaxed)}
(\texttt{ldrb}) rồi mới \texttt{exchange}. Phép đọc chỉ là gợi ý, \texttt{exchange} thắng mới là phép
acquire. Mọi vòng spin gọi \texttt{cpu\_relax()} giữa hai lần thử, tức là \texttt{pause} trên x86 và
\texttt{isb} trên AArch64. Lệnh \texttt{isb} được chọn thay cho lệnh \texttt{yield} của AArch64, vì trên
Apple \texttt{yield} gần như không làm gì (0,3~ns so với 12,1~ns của \texttt{isb}).
% nguồn: results/apple-m5/sleep_latency.csv: yield_insn trung vị 0,300 ns; isb trung vị 12,1 ns.

MCS~\cite{mellorcrummey1991mcs} giữ con trỏ \texttt{tail\_} tới cuối hàng đợi (\texttt{nullptr} là
rảnh). Để lấy lock, thread khởi tạo nút của mình bằng hai phép ghi relaxed, đổi \texttt{tail\_} sang nút
đó bằng \texttt{exchange} acq\_rel, rồi nếu hàng đợi khác rỗng thì nối nút vào \texttt{next} của nút
trước (release) và chờ cờ \texttt{locked} của chính nút mình (acquire). Để nhả lock, chủ lock đọc
\texttt{next}. Nếu chưa có thread kế tiếp, nó CAS \texttt{tail\_} về rỗng, còn nếu có thì đặt
\texttt{locked} của thread đó về \texttt{false} (release). Mỗi lần chuyển giao chỉ ghi một line của một
thread, theo thứ tự FIFO. Hai phép ghi relaxed an toàn vì nút chỉ trở nên với tới được qua hai phép
release theo sau. Nếu CAS
\texttt{tail\_} thất bại, một thread đã đổi \texttt{tail\_} nhưng chưa nối vào \texttt{next}, nên chủ
lock phải chờ nó. Một thread chỉ vào critical section khi nó thấy \texttt{tail\_} rỗng hoặc khi thread
trước nó đặt \texttt{locked} của nó về \texttt{false}, nên tại mỗi lúc chỉ có một chủ lock. Hàng đợi
theo thứ tự FIFO, nên không thread nào bị đói. Nút là biến \texttt{thread\_local} vì giao diện không có
tham số, nên một thread không giữ được hai lock MCS cùng lúc. Lock được giao cho đúng một thread, nên nếu
thread đó bị preempt (bị scheduler tước core để chạy thread khác) thì mọi thread phía sau phải chờ
scheduler.

\PCname{lock_tas_bo<B>} và \PCname{lock_ttas_bo<B>} gọi chính sách \texttt{B} sau mỗi lần thử thất bại.
TTAS chỉ lùi lại khi thua \texttt{exchange} (BackoffLock~\cite{herlihy2020art}). \PCname{sleep} theo
mẫu của đề, gọi \texttt{sleep\_for} với \texttt{delay} nhân đôi từ 100~ns tới \texttt{MAX\_DELAY}
$= 100\,000$~ns. \PCname{spin} quay ngẫu nhiên, với cửa sổ nhân đôi từ 4 tới 1024 lần \texttt{cpu\_relax}.
\PCname{yield} quay 64 lần rồi gọi \texttt{std::this\_thread::yield}, tức \texttt{sched\_yield}, để
nhường core cho thread khác. \texttt{sleep\_for} là system call. Trên
\PCname{apple-m5}, yêu cầu 100~ns ngủ trung vị 2\,670~ns. Lock thư viện (\texttt{std::mutex},
\texttt{pthread\_mutex\_t}, \texttt{os\_unfair\_lock}, \texttt{pthread\_spinlock\_t},
\texttt{omp\_lock\_t}) được bọc inline sau cùng giao diện.
% nguồn: results/apple-m5/sleep_latency.csv, sleep_for 100 ns: trung vị 2670 ns.
